\section{Proof-Theoretic Concepts and Basic TFL Metalogic}

So far I've been talking informally about valid \emph{arguments}. Something
like $\ephi,\epsi\therefore\ephi\eand\epsi$ outlines the argument, but we
want to be more precise in our discussion of various arguments. First, we
will replace the informal \therefore marker with a more formal meta-logical
operator \efollows and its negation \enotfollows.

\begin{definition}[Turnstiles]
  \label{def:turnstiles}
  Let $\ephi_1\dots\ephi_n$ and \epsi be TFL sentences. We can say
  $\ephi_1\dots\ephi_n\efollows\epsi$ iff there is a proof that ends in \epsi
  and has $\ephi_1\dots\ephi_n$ as its premises. We can say
  $\ephi_1\dots\ephi_n\enotfollows\epsi$ if there is no such proof.
\end{definition}

Now we can use the more precise turnstile notation to express a variety of
\emph{proof-theoretic} concepts. Many of these will be familiar from exploring
analogous \emph{truth-functional} concepts earlier.

First, we note a sentence which can be proven from nothing is a \emph{theorem}.

\begin{definition}[Theorem]
  \label{def:theorem}
  A sentence \ephi is a \emph{theorem} iff \efollows\ephi.
\end{definition}

Similarly, sentences that can be proven from each other are
\emph{interderivable}.

\begin{definition}[Interderivable]
  \label{def:interderivable}
  Two TFL sentences \ephi and \epsi are \emph{interderivable} iff
  \ephi\efollows\epsi and \epsi\efollows\ephi.
\end{definition}

\begin{definition}[Joint Consistency]
  \label{def:joint-consistency}
  TFL sentences $\ephi_1\dots\ephi_n$ are \emph{jointly inconsistent} iff
  $\ephi_1\dots\ephi_n\efollows\ebot$. Otherwise they are \emph{jointly
  consistent}.
\end{definition}

We can summarize the burdens of proof much as we did for truth tables.

\begin{center}
  \begin{tabular}{c c c}
                  & \textbf{yes} & \textbf{no} \\
    \hline
    theorem?      & one          & all         \\
    inconsistent? & one          & all         \\
    equivalent?   & two          & all         \\
  \end{tabular}
\end{center}

In this table we answer the question of how many proofs are needed to show the
property in question. The main takeaway is that, while proving the affirmative
is a straightfoward matter of logic, proving the negative requires metalogical
reasoning about \emph{all possible proofs}.

A special case of joint inconsistency is when there is only one premise. In
such a case, we can say that the sentence itself is \emph{inconsistent}.

\begin{definition}[Inconsistent Sentence]
  \label{def:consistent-sentence}
  A TFL sentence \ephi is \emph{inconsistent} iff $\ephi\efollows\ebot$.
\end{definition}

Every TFL sentence that isn't inconsistent or a theorem is \emph{contingent}.

\begin{definition}[Contingent Sentence]
  \label{def:contingent-sentence}
  A TFL sentence \ephi is \emph{(proof-theoretically) contingent} if \ephi is
  neither a theorem nor inconsistent.
\end{definition}

And we can make formal the notion of a \emph{valid argument} discussed
informally earlier as provability of the conclusion from the premises.

\begin{definition}[Valid Argument]
  \label{def:joint-consistency}
  The argument $\ephi_1\dots\ephi_n\therefore\epsi$ is a \emph{valid} iff
  $\ephi_1\dots\ephi_n\efollows\epsi$.
\end{definition}

From a metalogical perspective, there remains the question of the relationship
between the truth-functional and proof-theoretic approaches to logic. Meta-
logical proofs will be a major topic in Logic 2, but we will lay out the
broadstrokes before wrapping up our module on Truth-Functional Logic.

First, consider the how the definitions of the major logical concepts vary
between the two systems.

\begin{center}
  \begin{tabular}{c||c|c}
    \bf{concept} & \bf{truth-functional} & \bf{proof-theoretic} \\
    \hline
    Tautology / Theorem &
    always T &
    derived without premises \\
    Contradiction / Inconsistency &
    always F &
    can derive \ebot from it \\
    Contingent &
    sometimes T, sometimes F &
    not a theorem or inconsistency \\
    Equivalent / Interderivable &
    identical main columns &
    derive mutually \\
    Unsatisfiable / Inconsistent &
    no line where all T &
    can derive \ebot \\
    Satisfiable / Consistent &
    some line where all are T &
    cannot derive \ebot from them \\
    Valid argument &
    no lines w F $\ephi_i$ and T \epsi &
    can derive \epsi from $\ephi_1\dots\ephi_n$ \\
  \end{tabular}    
\end{center}

Two of the metalogical results that we will not prove here are \emph{soundness}
and \emph{consistency}. We will define them and them assume that they are true
as we wrap up our discussion of proofs for logical properties.

\begin{definition}[TFL Soundness]
  A proof system is \emph{sound} iff \ephi\efollows\epsi implies
  \ephi\eentails\epsi for all TFL sentences \ephi and \epsi.
\end{definition}

\begin{definition}[TFL Completeness]
  A proof system is \emph{complete} iff \ephi\eentails\epsi implies
  \ephi\efollows\epsi for all TFL sentences \ephi and \epsi.
\end{definition}

Knowing that a proof system is sound and complete means that we can demonstrate
a logical property using either proof-theoretic or truth-functional appoaches
and know that it holds for both. This means we can leverage the assymmetries in
the last table and approach the problem in whatever way is easiest. The
following table demonstrates such a strategy.

\begin{center}
  \begin{tabular}{c||c|c}
    \textbf{property} & \textbf{prove present} & \textbf{prove absent} \\
    \hline    
    Tautology, Theorem &
    derive the sentence &
    find a F line \\
    Contradiction &
    derive \ebot &
    find a T line \\
    Contingency &
    find T and F lines &
    prove \ephi or \enot\ephi \\
    Equivalence &
    mutually derive &
    find a line with truth mismatch \\
    Consistency &
    find a universally T line &
    derive \ebot \\
    Validity &
    derive \epsi from $\ephi_1\dots\ephi_n$ &
    find a line with  $\ephi_1\dots\ephi_n$ all T and \epsi F \\
  \end{tabular}    
\end{center}

